\section{Conclusion}

This paper tested whether a reinforcement learning agent can act as an
autonomous database administrator, moving shards in real time as the
geographic centre of a workload moves. Three things are established. Placement
changes 95th percentile latency on a real six-node cluster by 1372.5\,ms, so
the premise holds outside simulation. The room available to a better policy
grows with how fast the workload moves, from 0.44\,ms under slow drift to
57.20\,ms under fast drift, and at the fastest setting the reactive rule stops
beating fixed placement on the combined cost at all. And a Deep Q-Network
trained at that setting beats both baselines on held-out seeds, on two
independent training seeds, by moving one twentieth of the data the rule
moves.

The win is narrow, over an opponent that has itself stopped working, with a
policy that had not converged and that wastes most of its decisions on
inapplicable actions. It should be read as evidence that the room identified
by the sweep is reachable, not as a solved problem. The remaining distance to
the cost-free reference is larger than the distance already closed.

\section*{Disclosure of AI Use}

Claude, by Anthropic, was used throughout this project as a coding and writing
assistant. It helped implement the simulator, the baselines and the experiment
scripts, drafted and revised sections of this paper, and checked the
arithmetic in the result tables. All experiments were specified, dispatched
and reviewed by the author. Every number reported here was transcribed from a
recorded run directory, each recorded row names the run it came from, and the
cost arithmetic in Table~\ref{tab:agent} is reproduced by a script in the
project repository rather than computed by hand. The author is responsible for
the content and the claims made.
